\documentclass[11pt,a4paper]{article}
% ============================================================
% Shared preamble for MCA101 exam-preparation notes.
%
% libertinus and tcolorbox are installed in the user TeX tree
% (/home/elish4h/texmf) via:
%   tlmgr --usermode install --ignore-warning \
%     tcolorbox libertinus libertinus-type1 libertinus-otf \
%     libertinus-fonts libertinust1math
%
% The box styling matches the repository's notes-template.tex.
% ============================================================

\usepackage[margin=0.72in]{geometry}
\usepackage{libertinus}
\usepackage{microtype}
\usepackage[T1]{fontenc}
\usepackage{xcolor}
\usepackage{enumitem}
\usepackage{tabularx}
\usepackage{array}
\usepackage{booktabs}
\usepackage{longtable}
\usepackage{amsmath}
\usepackage{amssymb}
\usepackage{fancyhdr}
\usepackage{titlesec}
\usepackage[skins,breakable]{tcolorbox}
\usepackage{hyperref}
\usepackage{multicol}
\usepackage{parskip}

\definecolor{navy}{HTML}{183153}
\definecolor{blue}{HTML}{2563EB}
\definecolor{green}{HTML}{166534}
\definecolor{red}{HTML}{B91C1C}
\definecolor{gold}{HTML}{92400E}
\definecolor{lightblue}{HTML}{EFF6FF}
\definecolor{lightgreen}{HTML}{F0FDF4}
\definecolor{lightred}{HTML}{FEF2F2}
\definecolor{lightgold}{HTML}{FFFBEB}
\definecolor{grey}{HTML}{475569}

\hypersetup{colorlinks=true,linkcolor=blue,urlcolor=blue}

\titleformat{\section}{\Large\bfseries\color{navy}}{}{0pt}{}
\titleformat{\subsection}{\large\bfseries\color{navy}}{}{0pt}{}
\titleformat{\subsubsection}{\normalsize\bfseries\color{blue}}{}{0pt}{}

\setlist[itemize]{leftmargin=*,nosep}
\setlist[enumerate]{leftmargin=*,nosep}

% Allow tiny line-length overruns to be absorbed rather than overfull.
\setlength{\emergencystretch}{2em}

\tcbset{
  enhanced,
  breakable,
  colbacktitle=blue,
  coltitle=white,
  fonttitle=\bfseries,
  boxrule=0.5pt,
  arc=1.5pt,
  left=6pt,right=6pt,top=4pt,bottom=4pt,
  before skip=6pt,after skip=6pt,
}

\newtcolorbox{exambox}{colback=lightblue,colframe=blue,title=Exam Focus}
\newtcolorbox{rememberbox}{colback=lightgold,colframe=gold,title=Remember}
\newtcolorbox{dangerbox}{colback=lightred,colframe=red,title=Do Not Confuse}
\newtcolorbox{answerbox}{colback=lightgreen,colframe=green,title=Answer Skeleton}
\newtcolorbox{pyqbox}{colback=white,colframe=grey,title=Past-paper question}
\newtcolorbox{solbox}{colback=lightgreen,colframe=green,title=Solution}

\newcommand{\mkmarks}[1]{\hfill\textbf{[#1 marks]}}
\newcommand{\cmd}[1]{\texttt{#1}}
\newcommand{\F}[1]{\ensuremath{#1}}
\newcommand{\subject}[1]{\lhead{#1}\rhead{MCA101 Exam Notes}\cfoot{\thepage}}

\subject{MCA101 -- Module 4: Mathematical Logic}

\title{\textbf{Module 4: Mathematical Logic}\\}
\author{\large Propositional logic, equivalences, normal forms,\\predicates, quantifiers and rules of inference}
\date{}

\begin{document}
\maketitle

\begin{exambox}
\textbf{Scope.} Syllabus 4.1 to 4.3: mathematical logic (propositional and predicate logic, propositional equivalences, normal forms, predicates and quantifiers, nested quantifiers, rules of inference).

Logic appeared in four of the five papers in the set, always with an OR pair, and it is the module where the answer is a \emph{method} rather than a calculation. Learn the truth-table method and the algebraic method equally well, because both are asked.
\end{exambox}

\tableofcontents
\newpage

\section{Propositions and connectives}

A \textbf{proposition} is a declarative statement with exactly one truth value: true (T) or false (F). Questions, commands, and open sentences containing an unspecified variable are not propositions.

\begin{center}\small
\begin{tabularx}{\textwidth}{>{\bfseries}p{0.16\textwidth} p{0.18\textwidth} X X}
\toprule Connective & Notation & Meaning & Truth condition\\\midrule
Negation & $\neg p$ & not $p$ & True exactly when $p$ is false\\
Conjunction & $p\land q$ & $p$ and $q$ & True when both are true\\
Inclusive OR & $p\lor q$ & $p$ or $q$ or both & False only when both are false\\
Exclusive OR & $p\oplus q$ & exactly one of $p,q$ & True when their values differ\\
Implication & $p\to q$ & if $p$, then $q$ & False only for $p=T,q=F$\\
Biconditional & $p\leftrightarrow q$ & $p$ iff $q$ & True when their values agree\\
\end{tabularx}
\end{center}

For $n$ variables, a complete truth table has $2^n$ rows. Note the two key rewrites:
\[
p\to q\equiv\neg p\lor q,\qquad p\leftrightarrow q\equiv(p\to q)\land(q\to p).
\]

\begin{center}\small
\begin{tabular}{cc|cccccc}
\toprule $p$&$q$&$\neg p$&$p\land q$&$p\lor q$&$p\oplus q$&$p\to q$&$p\leftrightarrow q$\\\midrule
T&T&F&T&T&F&T&T\\
T&F&F&F&T&T&F&F\\
F&T&T&F&T&T&T&F\\
F&F&T&F&F&F&T&T\\\bottomrule
\end{tabular}
\end{center}

\section{Tautology, contradiction, contingency}

\begin{itemize}
\item A \textbf{tautology} is true on every row, e.g.\ $p\lor\neg p$.
\item A \textbf{contradiction} is false on every row, e.g.\ $p\land\neg p$.
\item A \textbf{contingency} is true on some rows and false on others, e.g.\ $p\land q$.
\end{itemize}
To classify a formula, construct its full truth table and inspect the final column.

\begin{pyqbox}
\textbf{Oct 2024 CAE-2 Q2A(a).} Check whether the statement $(p\lor\neg p)\lor\bigl((q\lor r)\lor\neg(q\lor r)\bigr)$ is a tautology or a contradiction.
\end{pyqbox}

\begin{solbox}
Simplify first. $(p\lor\neg p)$ is a tautology, so the whole formula is true regardless of the other part. Equally, $(q\lor r)\lor\neg(q\lor r)$ is a tautology. A disjunction of tautologies is a tautology.

Full truth table (8 rows for $p,q,r$):
\begin{center}\small
\begin{tabular}{ccc|ccccc}
\toprule $p$&$q$&$r$&$p\lor\neg p$&$q\lor r$&$\neg(q\lor r)$&$(q\lor r)\lor\neg(q\lor r)$&Full formula\\\midrule
T&T&T&T&T&F&T&T\\
T&T&F&T&T&F&T&T\\
T&F&T&T&T&F&T&T\\
T&F&F&T&F&T&T&T\\
F&T&T&T&T&F&T&T\\
F&T&F&T&T&F&T&T\\
F&F&T&T&T&F&T&T\\
F&F&F&T&F&T&T&T\\\bottomrule
\end{tabular}
\end{center}
The final column is T on every row, so the formula is a \textbf{tautology}.
\end{solbox}

\section{Logical equivalences}

Two formulas are logically equivalent, written $P\equiv Q$, when their truth values agree for every assignment.

\begin{center}\small
\begin{tabularx}{\textwidth}{>{\bfseries}p{0.22\textwidth} X X}
\toprule Law & Form 1 & Form 2\\\midrule
Identity & $p\land T$ & $p$; \quad $p\lor F=p$\\
Domination & $p\lor T$ & $T$; \quad $p\land F=F$\\
Idempotent & $p\lor p$ & $p$; \quad $p\land p=p$\\
Double negation & $\neg\neg p$ & $p$\\
Commutative & $p\lor q$ & $q\lor p$; similarly for $\land$\\
Associative & $(p\lor q)\lor r$ & $p\lor(q\lor r)$; similarly for $\land$\\
Distributive & $p\land(q\lor r)$ & $(p\land q)\lor(p\land r)$\\
De Morgan & $\neg(p\land q)$ & $\neg p\lor\neg q$\\
De Morgan & $\neg(p\lor q)$ & $\neg p\land\neg q$\\
Absorption & $p\lor(p\land q)$ & $p$; dual: $p\land(p\lor q)=p$\\
Implication & $p\to q$ & $\neg p\lor q$\\
Biconditional & $p\leftrightarrow q$ & $(p\to q)\land(q\to p)$\\
\end{tabularx}
\end{center}

\begin{dangerbox}
The distributive law also works in the other direction: $p\lor(q\land r)\equiv(p\lor q)\land(p\lor r)$. Students often remember only one version.
\end{dangerbox}

\subsection{Converse, inverse, contrapositive}
For $p\to q$: the \textbf{converse} is $q\to p$; the \textbf{inverse} is $\neg p\to\neg q$; the \textbf{contrapositive} is $\neg q\to\neg p$.

The original and the contrapositive are equivalent. The converse and the inverse are equivalent to each other, but are not generally equivalent to the original.

\begin{pyqbox}
\textbf{Nov 2024 Section A Q8.} Write the converse and the contrapositive of: "If a function is differentiable then it is continuous."
\end{pyqbox}

\begin{solbox}
Original: $p\to q$ where $p$ = "the function is differentiable", $q$ = "the function is continuous".

\textbf{Converse:} $q\to p$, that is, "If a function is continuous then it is differentiable." (False in general.)

\textbf{Contrapositive:} $\neg q\to\neg p$, that is, "If a function is not continuous then it is not differentiable." (True, since it is equivalent to the original.)

The inverse would be: "If a function is not differentiable then it is not continuous."
\end{solbox}

\begin{dangerbox}
For "if $p$, then $q$," do not swap the hypothesis and conclusion and assume the result is equivalent. Only the contrapositive is guaranteed equivalent.
\end{dangerbox}

\section{Proving equivalences}

Two methods are asked for, sometimes in the same paper as an OR pair.

\subsection{Truth-table method}
Build a truth table with a column for every intermediate sub-expression. If the two sides produce identical columns, they are equivalent.

\begin{pyqbox}
\textbf{Nov 2024 Q16(a)(i).} Using the truth table, prove the logical equivalence $\neg(p\leftrightarrow q)\equiv(p\land\neg q)\lor(q\land\neg p)$.
\end{pyqbox}

\begin{solbox}
\begin{center}\small
\begin{tabular}{cc|ccc|c|c}
\toprule $p$&$q$&$p\leftrightarrow q$&$p\land\neg q$&$q\land\neg p$&$(p\land\neg q)\lor(q\land\neg p)$&$\neg(p\leftrightarrow q)$\\\midrule
T&T&T&F&F&F&F\\
T&F&F&F&T&T&T\\
F&T&F&T&F&T&T\\
F&F&T&T&T&T&F\\\bottomrule
\end{tabular}
\end{center}
The final two columns are identical, so $\neg(p\leftrightarrow q)\equiv(p\land\neg q)\lor(q\land\neg p)$.
\end{solbox}

\begin{pyqbox}
\textbf{Nov 2024 Q16(b)(ii).} Using the truth table, prove the logical equivalence $p\land(q\lor r)\equiv(p\land q)\lor(p\land r)$.
\end{pyqbox}

\begin{solbox}
\begin{center}\small
\begin{tabular}{ccc|ccccc}
\toprule $p$&$q$&$r$&$q\lor r$&$p\land(q\lor r)$&$p\land q$&$p\land r$&$(p\land q)\lor(p\land r)$\\\midrule
T&T&T&T&T&T&T&T\\
T&T&F&T&T&T&F&T\\
T&F&T&T&T&F&T&T\\
T&F&F&F&F&F&F&F\\
F&T&T&T&F&F&F&F\\
F&T&F&T&F&F&F&F\\
F&F&T&T&F&F&F&F\\
F&F&F&F&F&F&F&F\\\bottomrule
\end{tabular}
\end{center}
The columns match, proving the distributive law. Note this is exactly the \emph{contradiction} test too: the formula is false only in the row $p=T,q=F,r=F$, where it is a contingency.
\end{solbox}

\subsection{Algebraic method (without a truth table)}
Start from one side, apply one named identity per step, and reach the other side.

\begin{pyqbox}
\textbf{Nov 2024 Q16(b)(i).} Without using a truth table, show that $p\leftrightarrow q\equiv(p\land q)\lor(\neg p\land\neg q)$.
\end{pyqbox}

\begin{solbox}
\[
\begin{aligned}
p\leftrightarrow q
&\equiv (p\to q)\land(q\to p) && \text{biconditional}\\
&\equiv (\neg p\lor q)\land(\neg q\lor p) && \text{implication}\\
&\equiv \bigl((\neg p\land\neg q)\lor(\neg p\land p)\bigr)\lor\bigl((q\land\neg q)\lor(q\land p)\bigr) && \text{distributive}\\
&\equiv (\neg p\land\neg q)\lor F\lor F\lor(q\land p) && \text{complement}\\
&\equiv (\neg p\land\neg q)\lor(p\land q) && \text{identity}\\
&\equiv (p\land q)\lor(\neg p\land\neg q) && \text{commutative}
\end{aligned}
\]
\end{solbox}

\begin{pyqbox}
\textbf{Aug 2025 Q3B(b).} Without using a truth table, prove that $\bigl[(p\lor q)\land\neg p\bigr]\to\neg q\equiv q\to p$.
\end{pyqbox}

\begin{solbox}
Work with the left side:
\[
\begin{aligned}
\bigl[(p\lor q)\land\neg p\bigr]\to\neg q
&\equiv \neg\bigl[(p\lor q)\land\neg p\bigr]\lor\neg q && \text{implication}\\
&\equiv \neg(p\lor q)\lor\neg\neg p\lor\neg q && \text{De Morgan, double negation}\\
&\equiv (\neg p\land\neg q)\lor p\lor\neg q && \text{De Morgan}\\
&\equiv p\lor(\neg p\land\neg q)\lor\neg q && \text{associative, commutative}\\
&\equiv p\lor\neg q && \text{absorption: }p\lor(\neg p\land\neg q)=p\\
&\equiv q\to p && \text{implication}
\end{aligned}
\]
\end{solbox}

\begin{exambox}
\textbf{Worked, tautology without a truth table.} Show $[(p\lor q)\land\neg p]\to q$ is a tautology.
\[
\begin{aligned}
\relax[(p\lor q)\land\neg p]\to q
&\equiv \neg[(p\lor q)\land\neg p]\lor q\\
&\equiv \neg(p\lor q)\lor p\lor q\\
&\equiv (\neg p\land\neg q)\lor p\lor q\\
&\equiv p\lor q\\
&\equiv T.
\end{aligned}
\]
Since the formula is equivalent to $T$, it is a tautology.
\end{exambox}

\section{Predicates and quantifiers}

A \textbf{predicate} $P(x)$ is a statement whose truth depends on the value assigned to $x$. A domain specifies the allowed values.

\begin{itemize}
\item $\forall x\,P(x)$: for every $x$ in the domain, $P(x)$ is true.
\item $\exists x\,P(x)$: there is at least one $x$ for which $P(x)$ is true.
\end{itemize}

\begin{rememberbox}
Negation rules: $\neg\forall x\,P(x)\equiv\exists x\,\neg P(x)$ and $\neg\exists x\,P(x)\equiv\forall x\,\neg P(x)$. For example, "not everyone passed" means "someone did not pass," not "nobody passed."
\end{rememberbox}

\subsection{Symbolizing statements}
\begin{center}\small
\begin{tabularx}{\textwidth}{>{\bfseries}p{0.40\textwidth} X}
\toprule Statement & Symbolic form\\\midrule
All men are giants & $\forall x(M(x)\to G(x))$\\
Some men are giants & $\exists x(M(x)\land G(x))$\\
No man is a giant & $\forall x(M(x)\to\neg G(x))$\\
Not every man is a giant & $\exists x(M(x)\land\neg G(x))$\\
All men are giants and all women are giants & $\forall xM(x)\to G(x)$ \emph{and} $\forall xW(x)\to G(x)$\\
\end{tabularx}
\end{center}

\begin{pyqbox}
\textbf{Oct 2025 Q4(a)(iii).} If $p$ means "he is intelligent" and $q$ means "he is strong", write the symbolic form of: "It is wrong that he is intelligent or strong."
\end{pyqbox}

\begin{solbox}
"It is wrong that" is a negation. So the statement is $\neg(p\lor q)$, which by De Morgan is $\neg p\land\neg q$, that is, "he is neither intelligent nor strong."
\end{solbox}

\begin{pyqbox}
\textbf{Nov 2024 Section A Q5.} Symbolize the statement "All men are giants".
\end{pyqbox}

\begin{solbox}
Let $M(x)$ be "x is a man" and $G(x)$ be "x is a giant".
\[
\forall x\,(M(x)\to G(x)).
\]
Note the implication, not a conjunction: a universal statement saying all men are giants does not assert that every object is both a man and a giant.
\end{solbox}

\subsection{Nested quantifiers}
$\forall x\exists y\,R(x,y)$ says each $x$ may have its own $y$; $\exists y\forall x\,R(x,y)$ says one fixed $y$ works for every $x$. In general these are not equivalent.

To disprove a universal statement, give a counterexample; to prove an existential statement, give a witness.

\begin{pyqbox}
\textbf{Oct 2025 Q4(b)(iii).} Represent the following in predicate calculus: (a) Some dragon is sleeping. (b) No dragon is sleeping. (c) Every dragon is sleeping. (d) Not every dragon is sleeping.
\end{pyqbox}

\begin{solbox}
Let $D(x)$ be "$x$ is a dragon" and $S(x)$ be "$x$ is sleeping".

(a) $\exists x\,(D(x)\land S(x))$.

(b) $\neg\exists x\,(D(x)\land S(x))$, equivalently $\forall x\,(D(x)\to\neg S(x))$.

(c) $\forall x\,(D(x)\to S(x))$.

(d) $\neg\forall x\,(D(x)\to S(x))$, equivalently $\exists x\,(D(x)\land\neg S(x))$.
\end{solbox}

\section{Rules of inference}

An argument is valid when its conclusion follows from its premises.

\begin{center}\small
\begin{tabularx}{\textwidth}{>{\bfseries}p{0.24\textwidth} X X}
\toprule Rule & Premises & Conclusion\\\midrule
Modus ponens & $p\to q$, $p$ & $q$\\
Modus tollens & $p\to q$, $\neg q$ & $\neg p$\\
Hypothetical syllogism & $p\to q$, $q\to r$ & $p\to r$\\
Disjunctive syllogism & $p\lor q$, $\neg p$ & $q$\\
Addition & $p$ & $p\lor q$\\
Simplification & $p\land q$ & $p$ (or $q$)\\
Conjunction & $p$, $q$ & $p\land q$\\
Resolution & $p\lor q$, $\neg p\lor r$ & $q\lor r$\\
\end{tabularx}
\end{center}

\begin{pyqbox}
\textbf{Nov 2024 Q16(a)(ii).} Explain the inference rules modus ponens and modus tollens.
\end{pyqbox}

\begin{solbox}
\textbf{Modus ponens:} from $p\to q$ and $p$, infer $q$. In words: if $p$ implies $q$ and $p$ is true, then $q$ is true.

\textbf{Modus tollens:} from $p\to q$ and $\neg q$, infer $\neg p$. In words: if $p$ implies $q$ and $q$ is false, then $p$ must be false.
\end{solbox}

\begin{pyqbox}
\textbf{Aug 2025 Q4A.} Show that the following set of premises is inconsistent: $P\to Q$, $P\to R$, $Q\to\neg R$, and $P$.
\end{pyqbox}

\begin{solbox}
Assume $P$. From $P\to Q$ and modus ponens, $Q$. From $P\to R$ and modus ponens, $R$. From $Q\to\neg R$ and modus ponens, $\neg R$. But $R$ and $\neg R$ cannot both be true.

Since the premises lead to a contradiction, the set is \textbf{inconsistent}.

Equivalently, from $P\to R$ we get $R$, and from $Q\to\neg R$ we get $Q\to\neg R$, which by modus tollens with $\neg R$ gives $\neg Q$. Combined with $Q$ this is also a contradiction.
\end{solbox}

\begin{pyqbox}
\textbf{Oct 2025 Q4(a)(ii).} Show that $\neg P$ is a valid conclusion from $P\to Q$ and $\neg Q$.
\end{pyqbox}

\begin{solbox}
This is modus tollens. Premise 1: $P\to Q$. Premise 2: $\neg Q$. Conclusion: $\neg P$.

\begin{center}\small
\begin{tabular}{lll}
\toprule & Statement & Reason\\\midrule
1 & $P\to Q$ & Premise\\
2 & $\neg Q$ & Premise\\
3 & $\neg P$ & 1, 2, modus tollens\\\bottomrule
\end{tabular}
\end{center}
\end{solbox}

\begin{pyqbox}
\textbf{Oct 2025 Q4(b)(i).} If A works hard, then either B or C will enjoy themselves. If B enjoys himself, then A will not work hard. If D enjoys himself, then C will not. Therefore, if A works hard, D will not enjoy himself. Show that these statements constitute a valid argument.
\end{pyqbox}

\begin{solbox}
Let $A$ = A works hard, $B$ = B enjoys, $C$ = C enjoys, $D$ = D enjoys.

Premises: $A\to(B\lor C)$; $B\to\neg A$; $D\to\neg C$.
Conclusion: $A\to\neg D$.

\begin{center}\small
\begin{tabular}{lll}
\toprule & Statement & Reason\\\midrule
1 & $A\to(B\lor C)$ & Premise\\
2 & $B\to\neg A$ & Premise\\
3 & $D\to\neg C$ & Premise\\
4 & $A$ & Assumption (conditional proof)\\
5 & $B\lor C$ & 1, 4, modus ponens\\
6 & $\neg B$ & 2 and 4, modus tollens\\
7 & $C$ & 5, 6, disjunctive syllogism\\
8 & $\neg C$ & 3 and 7, modus tollens\\
9 & $\neg D$ & 3 and 8, modus tollens\\
10 & $A\to\neg D$ & 4--9, conditional proof\\\bottomrule
\end{tabular}
\end{center}
Since line 10 follows from the premises, the argument is valid.
\end{solbox}

\section{Normal forms}

An \textbf{elementary product} (minterm) is an AND of literals; an \textbf{elementary sum} (maxterm) is an OR of literals. A literal is a variable or its negation.

\begin{itemize}
\item \textbf{DNF}: OR of AND terms, e.g.\ $(p\land q)\lor(\neg p\land r)$.
\item \textbf{CNF}: AND of OR terms, e.g.\ $(p\lor q)\land(\neg p\lor r)$.
\end{itemize}

To convert algebraically: remove $\to$ and $\leftrightarrow$, push negations inward using De Morgan's laws, then distribute. For canonical forms use the truth table: DNF takes the rows where the formula is true (minterms); CNF takes the rows where it is false (maxterms).

\begin{pyqbox}
\textbf{Oct 2025 Q4(a)(i).} Obtain the principal disjunctive normal forms of $\neg P\lor Q$.
\end{pyqbox}

\begin{solbox}
Build the truth table for $F=\neg P\lor Q$ over $P,Q$:

\begin{center}\small
\begin{tabular}{cc|cccc}
\toprule $P$&$Q$&$\neg P$&$Q$&F=$\neg P\lor Q$\\\midrule
T&T&F&T&T\\
T&F&F&F&F\\
F&T&T&T&T\\
F&F&T&F&T\\\bottomrule
\end{tabular}
\end{center}

$F$ is false only in the row $P=T,Q=F$. The principal DNF lists the minterm for every row where $F$ is true, with each variable appearing in its true (or negated) form:
\[
F=(P\land Q)\lor(\neg P\land Q)\lor(\neg P\land\neg Q).
\]
\end{solbox}

\begin{exambox}
\textbf{Worked, PDNF and CNF of $p\to q$.} The formula is false only when $p=T,q=F$.
\[
\text{PDNF}=(p\land q)\lor(\neg p\land q)\lor(\neg p\land\neg q),
\]
\[
\text{PCNF}=(\neg p\lor q).
\]
\end{exambox}

\begin{pyqbox}
\textbf{Oct 2025 Q4(a)(iv) and Nov 2024 Section A Q11.} What is conjunctive normal form?
\end{pyqbox}

\begin{solbox}
Conjunctive normal form (CNF) is a formula written as a conjunction (AND) of clauses, where each clause is a disjunction (OR) of literals. The canonical form, the principal conjunctive normal form, contains one maxterm for every row where the formula is false.

Example: $p\to(q\land r)$ has CNF $(\neg p\lor q)\land(\neg p\lor r)$.
\end{solbox}

\section{Answer skeletons}

\begin{answerbox}
\textbf{Truth table (3 to 9 marks).} List the variables; make all $2^n$ assignments; add a column for each intermediate connective; evaluate the target column; classify as tautology, contradiction or contingency and state the result.
\end{answerbox}

\begin{answerbox}
\textbf{Equivalence proof (5 to 9 marks).} Write the expression; remove implications; apply the named equivalence law one step at a time; reach the target expression. Alternatively, show identical truth-table columns. Name each law to make the reasoning clear.
\end{answerbox}

\begin{answerbox}
\textbf{Inference or validity (5 to 9 marks).} Number the premises; identify the rule applied at each line; show the conclusion as the last line. For a conditional conclusion, state the assumption and use conditional proof.
\end{answerbox}

\begin{answerbox}
\textbf{Quantifier negation (3 marks).} Negate the quantifier, switch $\forall\leftrightarrow\exists$, then negate the predicate.
\end{answerbox}

\begin{answerbox}
\textbf{Normal form (3 to 9 marks).} Write the truth table, identify the rows where the formula is true (for DNF) or false (for CNF), then write one minterm or maxterm per row with all variables present.
\end{answerbox}

\section{Self-test}

\begin{enumerate}
\item Define proposition, predicate, tautology, contradiction and contingency.
\item State the truth condition for implication and biconditional.
\item State De Morgan's laws and the negation rules for quantifiers.
\item Define DNF and CNF, with one example of each.
\item Find the converse, inverse and contrapositive of "If a number is divisible by 4, then it is even."
\item Negate: "Every student solved some problem." Express the result with quantifiers.
\item Identify the inference rule: $p\to q$, $q\to r$, therefore $p\to r$.
\item Prove $p\to(q\to p)$ is a tautology without a truth table.
\item Prove by truth table that $p\land(q\lor r)\equiv(p\land q)\lor(p\land r)$.
\item Obtain the principal disjunctive normal form of $\neg p\lor q$.
\item Show that the premises $p\to q$, $q\to r$, $r\to\neg p$ are inconsistent.
\item Represent in predicate calculus: every student passed at least one exam.
\end{enumerate}

\subsection*{Answers}
\begin{enumerate}
\item Proposition: a declarative statement with one truth value. Predicate: a statement whose truth depends on an unspecified variable. Tautology: true on every row. Contradiction: false on every row. Contingency: true on some rows and false on others.
\item $p\to q$ is false only when $p=T,q=F$; $p\leftrightarrow q$ is true exactly when $p$ and $q$ agree.
\item $\neg(p\land q)\equiv\neg p\lor\neg q$; $\neg(p\lor q)\equiv\neg p\land\neg q$; $\neg\forall xP\equiv\exists x\neg P$; $\neg\exists xP\equiv\forall x\neg P$.
\item DNF is an OR of AND terms, e.g.\ $(p\land q)\lor(\neg p\land r)$. CNF is an AND of OR terms, e.g.\ $(p\lor q)\land(\neg p\lor r)$.
\item Converse: "If a number is even then it is divisible by 4." Inverse: "If a number is not divisible by 4 then it is not even." Contrapositive: "If a number is not even then it is not divisible by 4." The original and the contrapositive are equivalent.
\item $\neg\bigl[\forall x\,\exists y\,S(x,y)\bigr]\equiv\exists x\,\forall y\,\neg S(x,y)$: there is a student who solved no exam.
\item Hypothetical syllogism.
\item $p\to(q\to p)\equiv\neg p\lor(\neg q\lor p)\equiv p\lor\neg p\lor\neg q\equiv T$.
\item Eight-row truth table; the columns for $p\land(q\lor r)$ and $(p\land q)\lor(p\land r)$ are identical.
\item $(p\land q)\lor(\neg p\land q)\lor(\neg p\land\neg q)$.
\item From $p\to q$ and $q\to r$ and $r\to\neg p$, assume $p$ to get $q$, then $r$, then $\neg p$, a contradiction. Hence inconsistent.
\item $\forall x\,\exists y\,S(x,y)$.
\end{enumerate}

\section{Last-minute page}

\begin{multicols}{2}
\textbf{CONNECTIVES}\\
$p\to q\equiv\neg p\lor q$\\
$p\leftrightarrow q\equiv(p\to q)\land(q\to p)$\\
$n$ variables $\Rightarrow 2^n$ rows\\[4pt]
\textbf{LAWS}\\
$\neg(p\land q)\equiv\neg p\lor\neg q$\\
$\neg(p\lor q)\equiv\neg p\land\neg q$\\
$p\land(q\lor r)\equiv(p\land q)\lor(p\land r)$\\
$p\lor(q\land r)\equiv(p\lor q)\land(p\lor r)$\\
$p\lor(p\land q)\equiv p$\\[4pt]
\textbf{QUANTIFIERS}\\
$\neg\forall xP\equiv\exists x\neg P$\\
$\neg\exists xP\equiv\forall x\neg P$\\
$\forall x\exists yR \ne \exists y\forall xR$\\[4pt]
\textbf{INFERENCE}\\
Modus ponens: $p\to q,\,p\Rightarrow q$\\
Modus tollens: $p\to q,\,\neg q\Rightarrow\neg p$\\
Hyp. syllogism, disj. syllogism, resolution\\[4pt]
\textbf{NORMAL FORMS}\\
DNF: OR of AND terms\\
CNF: AND of OR terms\\
PDNF: minterms of the TRUE rows\\
PCNF: maxterms of the FALSE rows\\[4pt]
\textbf{DO NOT CONFUSE}\\
Only the contrapositive matches the original.\\
$\forall$ vs $\exists$ changes meaning on negation.\\
$P$ alone is not a proposition.
\end{multicols}

\vfill
\begin{center}\textit{Name the law at each step. "Transforming" without naming the law earns half marks in a 9-mark proof.}\end{center}

\end{document}
